
\documentclass{article}
\usepackage{colortbl}
\usepackage{makecell}
\usepackage{multirow}
\usepackage{supertabular}

\begin{document}

\newcounter{utterance}

\centering \large Interaction Transcript for game `eqbench', experiment `eq\_bench', episode 33 with reflexion{-}llama{-}3.1{-}8b{-}instruct.
\vspace{24pt}

{ \footnotesize  \setcounter{utterance}{1}
\setlength{\tabcolsep}{0pt}
\begin{supertabular}{c@{$\;$}|p{.15\linewidth}@{}p{.15\linewidth}p{.15\linewidth}p{.15\linewidth}p{.15\linewidth}p{.15\linewidth}}
    \# & \multicolumn{2}{c}{Player} && \multicolumn{2}{c}{Game Master} \\
    \hline

    \theutterance \stepcounter{utterance}  
    & & & \multicolumn{4}{p{0.6\linewidth}}{
        \cellcolor[rgb]{0.9,0.9,0.9}{
            \makecell[{{p{\linewidth}}}]{
                \texttt{\tiny{[P1$\langle$GM]}}
                \texttt{You have attempted this same task before and did not succeed. These are the plans you wrote after those attempts.} \\
\texttt{Trial \#1: New plan: } \\
\\ 
\texttt{Upon re{-}reading the dialogue, I realized that I focused too much on the surface{-}level conflict between Brenda and Tessa, but didn't delve deeper into the underlying emotions and motivations. I should have paid more attention to the power dynamics and the emotional nuances of the conversation. } \\
\\ 
\texttt{In particular, I should have noticed that Brenda's responses are more defensive and controlling, while Tessa's responses are more assertive and emotional. This suggests that Brenda is feeling a sense of fear and frustration, not just because of the conflict itself, but because she's losing control of the situation and her daughter is pushing back against her.} \\
\\ 
\texttt{On the other hand, Tessa's responses suggest that she's feeling a sense of helplessness and frustration, not just because of her mother's behavior, but because she's struggling to be heard and understood. However, Tessa's final statement, "Maybe I need to make my own mistakes, Mom," suggests that she's also feeling a sense of acceptance and determination to assert her independence.} \\
\\ 
\texttt{In my next attempt, I will focus on analyzing the emotional nuances of the conversation, paying attention to the power dynamics, and identifying the underlying emotions and motivations of the characters. I will also consider the context and the relationships between the characters to provide a more accurate and empathetic response.} \\
\\ 
\texttt{Trial \#2: New plan: } \\
\\ 
\texttt{Upon re{-}reading the dialogue, I realized that I focused too much on the surface{-}level conflict between Brenda and Tessa, but didn't delve deeper into the underlying emotions and motivations. I should have paid more attention to the power dynamics and the emotional nuances of the conversation. } \\
\\ 
\texttt{In particular, I should have noticed that Brenda's responses are more defensive and controlling, while Tessa's responses are more assertive and emotional. This suggests that Brenda is feeling a sense of fear and frustration, not just because of the conflict itself, but because she's losing control of the situation and her daughter is pushing back against her.} \\
\\ 
\texttt{However, I also noticed that I didn't give enough weight to the fact that Brenda is trying to protect her daughter from making the same mistakes she did. This suggests that Brenda's fear and frustration are also motivated by a desire to prevent harm to her daughter, rather than just a desire to control the situation.} \\
\\ 
\texttt{In my next attempt, I will focus on analyzing the emotional nuances of the conversation, paying attention to the power dynamics, and identifying the underlying emotions and motivations of the characters. I will also consider the context and the relationships between the characters to provide a more accurate and empathetic response.} \\
\\ 
\texttt{To improve my analysis, I will ask myself the following questions:} \\
\\ 
\texttt{* What are the underlying emotions and motivations of each character?} \\
\texttt{* How do the power dynamics between the characters influence their behavior and interactions?} \\
\texttt{* What are the emotional nuances of the conversation, and how do they reveal the underlying emotions and motivations of the characters?} \\
\texttt{* How does the context and the relationships between the characters influence their behavior and interactions?} \\
\\ 
\texttt{By asking these questions and considering the emotional nuances of the conversation, I will be able to provide a more accurate and empathetic response that takes into account the complexities of human emotions and relationships.} \\
\\ 
\texttt{Trial \#3: New plan: } \\
\\ 
\texttt{Upon re{-}reading the dialogue, I realized that I focused too much on the surface{-}level conflict between Brenda and Tessa, but didn't delve deeper into the underlying emotions and motivations. I should have paid more attention to the power dynamics and the emotional nuances of the conversation. } \\
\\ 
\texttt{In particular, I should have noticed that Brenda's responses are more defensive and controlling, while Tessa's responses are more assertive and emotional. This suggests that Brenda is feeling a sense of fear and frustration, not just because of the conflict itself, but because she's losing control of the situation and her daughter is pushing back against her.} \\
\\ 
\texttt{However, I also noticed that I didn't give enough weight to the fact that Brenda is trying to protect her daughter from making the same mistakes she did. This suggests that Brenda's fear and frustration are also motivated by a desire to prevent harm to her daughter, rather than just a desire to control the situation.} \\
\\ 
\texttt{In my next attempt, I will focus on analyzing the emotional nuances of the conversation, paying attention to the power dynamics, and identifying the underlying emotions and motivations of the characters. I will also consider the context and the relationships between the characters to provide a more accurate and empathetic response.} \\
\\ 
\texttt{To improve my analysis, I will ask myself the following questions:} \\
\\ 
\texttt{* What are the underlying emotions and motivations of each character?} \\
\texttt{* How do the power dynamics between the characters influence their behavior and interactions?} \\
\texttt{* What are the emotional nuances of the conversation, and how do they reveal the underlying emotions and motivations of the characters?} \\
\texttt{* How does the context and the relationships between the characters influence their behavior and interactions?} \\
\\ 
\texttt{By asking these questions and considering the emotional nuances of the conversation, I will be able to provide a more accurate and empathetic response that takes into account the complexities of human emotions and relationships.} \\
\\ 
\texttt{To improve my scoring, I will also consider the following:} \\
\\ 
\texttt{* Fear: 9 (Brenda's defensive and controlling behavior suggests a deep{-}seated fear of losing control and her daughter making mistakes)} \\
\texttt{* Helplessness: 3 (Tessa's assertive and emotional responses suggest that she's not feeling helpless, but rather empowered to express her feelings and needs)} \\
\texttt{* Frustration: 8 (Both Brenda and Tessa are feeling frustrated, but in different ways {-} Brenda is frustrated by Tessa's behavior, while Tessa is frustrated by her mother's controlling behavior)} \\
\texttt{* Acceptance: 7 (Tessa's final statement suggests that she's accepting of her mother's limitations and is determined to make her own mistakes and learn from them)} \\
\\ 
\\ 
\\ 
\texttt{Your task is to predict the likely emotional responses of a character in this dialogue:} \\
\\ 
\texttt{Brenda: So, you thought it was a good idea to turn our home into a nightclub?} \\
\texttt{Tessa: It wasn't a nightclub, Mom. Just a few friends hanging out.} \\
\texttt{Brenda: Hanging out? The neighbors called the police, Tessa!} \\
\texttt{Tessa: Maybe if you weren't so uptight, they wouldn't have to.} \\
\texttt{Brenda: This is not about me being uptight. This is about you respecting our home.} \\
\texttt{Tessa: Our home? Or your home?} \\
\texttt{Brenda: Don't start that again, Tessa.} \\
\texttt{Tessa: What, the truth? That you'd rather have a perfect little doll than a daughter?} \\
\texttt{Brenda: That's not fair. I just want you to be responsible.} \\
\texttt{Tessa: No, you want me to be you.} \\
\texttt{Brenda: And what's so wrong with that? I've worked hard to give us a good life.} \\
\texttt{Tessa: A good life? Or a perfect life?} \\
\texttt{Brenda: What's the difference?} \\
\texttt{Tessa: The difference is that in a good life, people make mistakes. They mess up, they learn, they grow. In your perfect life, there's no room for error. No room for me.} \\
\texttt{Brenda: That's not true, Tessa. I just... I don't want you to make the same mistakes I did.} \\
\texttt{Tessa: Maybe I need to make my own mistakes, Mom. } \\
\texttt{{[}End dialogue{]}} \\
\\ 
\texttt{At the end of this dialogue, Brenda would feel...} \\
\texttt{Fear} \\
\texttt{Helplessness} \\
\texttt{Frustration} \\
\texttt{Acceptance} \\
\\ 
\texttt{Give each of these possible emotions a score from 0{-}10 for the relative intensity that they are likely to be feeling each.} \\
\\ 
\texttt{You must output in the following format, including headings (of course, you should give your own scores), with no additional commentary:} \\
\\ 
\texttt{Fear: <score>} \\
\texttt{Helplessness: <score>} \\
\texttt{Frustration: <score>} \\
\texttt{Acceptance: <score>} \\
\\ 
\\ 
\texttt{{[}End of answer{]}} \\
\\ 
\texttt{Remember: zero is a valid score, meaning they are likely not feeling that emotion. You must score at least one emotion > 0.} \\
\\ 
\texttt{Your answer:} \\
            }
        }
    }
     \\ \\

    \theutterance \stepcounter{utterance}  
    & \multicolumn{4}{p{0.6\linewidth}}{
        \cellcolor[rgb]{0.561,0.737,0.561}{
            \makecell[{{p{\linewidth}}}]{
                \texttt{\tiny{[P1$\rangle$GM]}}
                \texttt{Fear: 9} \\
\texttt{Helplessness: 0} \\
\texttt{Frustration: 8} \\
\texttt{Acceptance: 6} \\
            }
        }
    }
    & & \\ \\

    \theutterance \stepcounter{utterance}  
    & & & \multicolumn{2}{p{0.3\linewidth}}{
        \cellcolor[rgb]{0.9,0.9,0.9}{
            \makecell[{{p{\linewidth}}}]{
                \texttt{\tiny{[GM$|$GM]}}
                \texttt{Fear: 9} \\
\texttt{Helplessness: 0} \\
\texttt{Frustration: 8} \\
\texttt{Acceptance: 6} \\
            }
        }
    }
    & & \\ \\

    \theutterance \stepcounter{utterance}  
    & & & \multicolumn{2}{p{0.3\linewidth}}{
        \cellcolor[rgb]{0.9,0.9,0.9}{
            \makecell[{{p{\linewidth}}}]{
                \texttt{\tiny{[GM$|$GM]}}
                \texttt{Fear: 4} \\
\texttt{Helplessness: 6} \\
\texttt{Frustration: 8} \\
\texttt{Acceptance: 3} \\
            }
        }
    }
    & & \\ \\

    \theutterance \stepcounter{utterance}  
    & & & \multicolumn{2}{p{0.3\linewidth}}{
        \cellcolor[rgb]{0.9,0.9,0.9}{
            \makecell[{{p{\linewidth}}}]{
                \texttt{\tiny{[GM$|$GM]}}
                \texttt{game\_result = LOSE} \\
            }
        }
    }
    & & \\ \\

    \theutterance \stepcounter{utterance}  
    & & \multicolumn{4}{p{0.6\linewidth}}{
        \cellcolor[rgb]{0.95,0.95,0.95}{
            \makecell[{{p{\linewidth}}}]{
                \texttt{\tiny{[P1$\leftrightarrow$P]}}
                \texttt{Post{-}episode Reflexion (not sent to GM):} \\
\texttt{New plan: } \\
\\ 
\texttt{Upon re{-}reading the dialogue, I realized that I focused too much on the surface{-}level conflict between Brenda and Tessa, but didn't delve deeper into the underlying emotions and motivations. I should have paid more attention to the power dynamics and the emotional nuances of the conversation. } \\
\\ 
\texttt{In particular, I should have noticed that Brenda's responses are more defensive and controlling, while Tessa's responses are more assertive and emotional. This suggests that Brenda is feeling a sense of fear and frustration, not just because of the conflict itself, but because she's losing control of the situation and her daughter is pushing back against her.} \\
\\ 
\texttt{However, I also noticed that I didn't give enough weight to the fact that Brenda is trying to protect her daughter from making the same mistakes she did. This suggests that Brenda's fear and frustration are also motivated by a desire to prevent harm to her daughter, rather than just a desire to control the situation.} \\
\\ 
\texttt{In my next attempt, I will focus on analyzing the emotional nuances of the conversation, paying attention to the power dynamics, and identifying the underlying emotions and motivations of the characters. I will also consider the context and the relationships between the characters to provide a more accurate and empathetic response.} \\
\\ 
\texttt{To improve my analysis, I will ask myself the following questions:} \\
\\ 
\texttt{* What are the underlying emotions and motivations of each character?} \\
\texttt{* How do the power dynamics between the characters influence their behavior and interactions?} \\
\texttt{* What are the emotional nuances of the conversation, and how do they reveal the underlying emotions and motivations of the characters?} \\
\texttt{* How does the context and the relationships between the characters influence their behavior and interactions?} \\
\\ 
\texttt{By asking these questions and considering the emotional nuances of the conversation, I will be able to provide a more accurate and empathetic response that takes into account the complexities of human emotions and relationships.} \\
\\ 
\texttt{To improve my scoring, I will also consider the following:} \\
\\ 
\texttt{* Fear: 9 (Brenda's defensive and controlling behavior suggests a deep{-}seated fear of losing control and her daughter making mistakes)} \\
\texttt{* Helplessness: 3 (Tessa's assertive and emotional responses suggest that she's not feeling helpless, but rather empowered to express her feelings and needs)} \\
\texttt{* Frustration: 8 (Both Brenda and Tessa are feeling frustrated, but in different ways {-} Brenda is frustrated by Tessa's behavior, while Tessa is frustrated by her mother's controlling behavior)} \\
\texttt{* Acceptance: 7 (Tessa's final statement suggests that she's accepting of her mother's limitations and is determined to make her own mistakes and learn from them)} \\
            }
        }
    }
    & \\ \\

\end{supertabular}
}

\end{document}
